% ----------------------------------------------------------
% Introdução
% ----------------------------------------------------------
\chapter{Introdução}

Na rotina do comércio artesanal e das confeitarias, a transição entre a gestão intuitiva e uma administração estruturada costuma ser comprometida pela descentralização de informações, registros informais e pela dependência da memória para o controle de encomendas e estoque. Conforme o negócio expande, o uso de métodos manuais resulta em falhas operacionais, retrabalho e perda de produtividade \cite{sobrenomeAno}.

A aplicação da tecnologia surge como uma alternativa viável para a organização de pequenos empreendimentos, otimizando processos e minimizando desperdícios. No setor de confeitaria, o gerenciamento de prazos, produtos personalizados e insumos demanda um controle rigoroso e contínuo \cite{sobrenomeAno}. 

Diante desse cenário, este Projeto de Conclusão de Curso apresenta o desenvolvimento do Sistema Web Administrativo para a confeitaria Bom Gosto. A solução proposta visa integrar o recebimento de pedidos à rotina da cozinha e ao controle de estoque, automatizando o fluxo operacional da empresa. O propósito central consiste em substituir os processos manuais por uma ferramenta digital padronizada que proporcione eficiência operacional e assertividade à gestão.

\section{Caracterização do Projeto}
O projeto propõe o desenvolvimento de um Sistema Web Administrativo para a confeitaria Bom Gosto,
com o objetivo de atender às suas necessidades operacionais e contribuir para a organização de seus
processos internos.

O sistema será voltado à administração da confeitaria, contemplando o gerenciamento das informações 
relacionadas aos pedidos, produtos e estoque. Dessa forma, a solução busca substituir métodos manuais
por um método mais ágil e preciso, proporcionando maior organização e eficiência às atividades administrativas.

\section{Objetivos}

\subsection{Objetivo Geral} 
Em pequenos negócios artesanais, a rotina corrida muitas vezes dificulta o acompanhamento organizado das finanças e dos pedidos. A tecnologia surge como um recurso indispensável para simplificar tarefas, eliminar controles manuais repetitivos e trazer clareza para o funcionamento diário da empresa. Desta forma, o objetivo é desenvolver uma plataforma web sob medida para a confeitaria Bom Gosto, capaz de centralizar as tarefas administrativas e substituir as anotações em papel por um fluxo digital prático e rápido. 

\subsection{Objetivos Específicos}
\begin{itemize}
    \item possibilitar o registro e o acompanhamento detalhado de novos pedidos e o controle desde solicitação do cliente até a entrega final;
    \item permitir o acompanhamento dos pedidos registrados no sistema, proporcionando maior organização das encomendas;
    \item possibilitar o gerenciamento dos produtos produzidos pela confeitaria;
    \item permitir o acompanhamento das informações relacionadas ao estoque da loja;
    \item integrar essas funcionalidades em um sistema automatizado que proporcione acesso às informações necessárias, modernizando a gestão operacional e reduzindo a entrada manual de dados.
\end{itemize}

\section{Justificativa}
Substituir as anotações manuais por uma solução digital reduz drasticamente o risco de perda de informações e erros de registro no dia a dia da confeitaria. Com um sistema organizando o fluxo de trabalho, os processos administrativos tornam-se padronizados e mais rápidos, evitando retrabalho na cozinha e no atendimento.

Essa organização traz mais estabilidade para a gestão da confeitaria Bom Gosto. Ao diminuir o tempo gasto com tarefas repetitivas de controle, a proprietária ganha mais margem para focar na qualidade dos produtos, no relacionamento com os clientes e no crescimento sustentável do negócio.

\section{Escopo do Projeto}

O projeto compreende o desenvolvimento de um Sistema Web Administrativo destinado à organização dos processos internos da confeitaria Bom Gosto.

Compreendem o escopo do projeto:
\begin{itemize}
    \item o registro e acompanhamento de pedidos;
    \item o gerenciamento dos produtos produzidos pela confeitaria;
    \item o controle das informações relacionadas ao estoque;
    \item a organização das informações administrativas em um sistema web.
\end{itemize}

Não fazem parte do escopo do projeto: 
\begin{itemize}
    \item o desenvolvimento de um website público de divulgação da confeitaria;
    \item a implementação de uma área de avaliações ou feedback público;
    \item a implementação de sistema de vendas ou ambiente transacional online.
\end{itemize}

\section{Estrutura do Documento}
Este trabalho está estruturado em seis capítulos, dispostos de forma a guiar o leitor desde a contextualização do problema até a concretização do software:

\begin{itemize}
    \item Capítulo 1 -- Introdução: Apresenta o projeto, contextualizando o leitor quanto aos seus objetivos, justificativa, escopo e estrutura geral, com foco no valor de negócio e nos aspectos funcionais da solução;
    \item Capítulo 2 -- Contexto do Projeto: Detalha o cenário real de aplicação, o público-alvo, as necessidades identificadas e as regras de negócio do parceiro atendido, fundamentando o problema antes da abordagem técnica;
    \item Capítulo 3 -- Análise do Sistema: Mapeia os requisitos, atores, restrições e modelos conceituais, estabelecendo as funcionalidades sem ainda aprofundar nas especificidades técnicas de implementação;
    \item Capítulo 4 -- Projeto do Sistema: Converte o modelo conceitual em uma especificação formal, definindo a arquitetura, componentes, modelagem de dados, interfaces e tecnologias adotadas;
    \item Capítulo 5 -- Implementação: Descreve a construção efetiva do software, abordando o ambiente de desenvolvimento, a integração dos módulos, os testes realizados, os resultados obtidos e as limitações encontradas;
    \item Capítulo 6 -- Conclusão: Apresenta as considerações finais, a avaliação dos resultados frente aos objetivos propostos, as dificuldades enfrentadas e as propostas para trabalhos futuros.
\end{itemize}

